\section{Estimating the limiting metric modulus from earlier data}
\label{sec:metric-modulus}
The polarized Einstein--scalar benchmark distinguishes the
additive metric datum from the derivative and shape fields.
Its production initial profile is a finite Fourier sum.
This permits an independent continuum computation of the
limiting modulus, and clarifies a more general stability
question for finite observations. The logarithmic term in
the field intercept is the established Kasner scattering
correction \cite[Theorem~1.3 and Section~1.4]{li}.
The result below concerns its contribution to the nonlinear
pointwise metric modulus after inversion of two field
snapshots.

Use the positive-quadrature convention with
$\tau=-\log t$ and scalar $\phi=\sqrt2\varphi$, where
$\varphi$ is canonically normalized at $8\pi G=1$.
The orbit field $P$ and scalar $\phi$ each satisfy
$u_{tt}+t^{-1}u_t-u_{\theta\theta}=0$.
The metric coefficient transverse to the two
Killing directions is $t^{-1/2}e^{-\lambda/2}$.
With $U=(P,\phi)$ one has
$\lambda_\tau=|U_\tau|^2+e^{-2\tau}|U_\theta|^2$.
For a fixed angular origin define
$v=\lim_{\tau\to\infty}U_\tau$,
$U_\infty=\lim_{\tau\to\infty}(U-\tau v)$,
and $c_\lambda=\lim_{\tau\to\infty}(\lambda(\tau,0)-\tau|v(0)|^2)$.
The origin can be replaced by any fixed angle.

At $t=a$, with $\tau_a=-\log a$ and one measured metric anchor,
\begin{align}
 c_\lambda
 &=\lambda(\tau_a,0)-\tau_a|v(0)|^2
   +\int_{\tau_a}^\infty
     \left(|U_\tau|^2-|v|^2+e^{-2\tau}|U_\theta|^2\right)(\tau,0)
     \,d\tau \nonumber\\
 &=\lambda(\tau_a,0)+\tau_a|v(0)|^2
     +2v(0)\cdot\big(U_\infty(0)-U(\tau_a,0)\big)
     +\int_{\tau_a}^\infty
        \left(|U_\tau-v|^2+e^{-2\tau}|U_\theta|^2\right)(\tau,0)
        \,d\tau .
 \label{eq:modulus-identity}
\end{align}
The second equality follows by integrating the part
linear in $U_\tau-v$. It displays the dependence on
the intercept as well as the slope. Both formulas
use one finite metric anchor, whose additive change
adds the same constant to $c_\lambda$.

\subsection{A sharp logarithmic threshold for the missing modulus}
For $\alpha\in\mathbb R$, let $\mathcal H^2_\alpha$
denote the space with Fourier coefficient weight
$\langle k\rangle^2[\log(2+|k|)]^\alpha$ in $\ell^2$.
Consider two full orbit-and-scalar snapshots at
$t=a,b$, with rational $b-a$ in units of $\pi$
and no zero modal transfer at that pair.
The excluded reference times form a locally finite
set at each such gap, by the scalar phase estimate.

\begin{theorem}[Sharp threshold for the pointwise modulus]
\label{th:metric-modulus}
With one metric anchor at the fixed angle and time $a$,
the limiting metric modulus in the polarized
Einstein--scalar system is locally Lipschitz in the
two snapshots with the $\mathcal H^2_\alpha$ norm if
$\alpha>3/2$. In particular it is locally Lipschitz
in $H^r$ for every $r>2$.
Near a homogeneous datum with nonzero limiting
velocity, a local Lipschitz estimate fails on smooth
periodic, constraint-compatible data for
$\alpha\le3/2$, and hence in the unweighted
Sobolev space $H^2$ or lower.
\end{theorem}

\begin{proof}
For either field, write the nonzero modes as
$A_kJ_0(|k|t)+B_kY_0(|k|t)$.
At the stated observation pair the inverse Bessel
matrix gives an upper bound
\[
 |A_k|+|B_k|\le C\langle k\rangle^{3/2}
                 \big(|\widehat U(a,k)|+|\widehat U(b,k)|\big).
\]
Indeed the Hankel expansion for $\beta=1/4$
gives order $k^{-1}$ for the energy-normalized
transfer on the resonant subsequence and a uniform
lower bound away from that subsequence. The finite
nonzero transfers and the scalar zero mode pose
no further difficulty. Thus the snapshots bound
the coefficient norm with weight
$\langle k\rangle^{1/2}[\log(2+|k|)]^\alpha$.
The limiting mode formulas are
\[
 v_k=-\frac{2B_k}{\pi},\qquad
 (U_\infty)_k=A_k+
      \frac{2B_k}{\pi}\big(\log(|k|/2)+\gamma_{\mathrm E}\big).
\]
Their absolute Fourier sums are bounded for
$\alpha>3/2$, by Cauchy--Schwarz.

We show that the integrals in the second line of
\eqref{eq:modulus-identity} are bounded quadratic
forms in the same data norm. At fixed $\theta$,
in the measure $dt/t$, consider
$G=U_\tau-v$ and $E=tU_\theta$.
For $z=|k|t\le1$, Bessel estimates give
\[
 |G_k|\le C(|A_k|+|B_k|)z^2(1+|\log z|),\qquad
 |E_k|\le C(|A_k|+|B_k|)z(1+|\log z|).
\]
The $L^2(dt/t)$ norm of each bound is at most a
constant times the coefficient magnitude,
uniformly in $k$. Their sum is finite in the stated
coefficient norm.

For $z\ge1$, the leading terms of both
$U_\tau$ and $E$ are sums of
$\sqrt{|k|t}\,c_k^\pm e^{\pm i|k|t}$
with $|c_k^\pm|\le C(|A_k|+|B_k|)$.
After removing the common $\sqrt t$, untruncated
sums are bounded in $L^2(0,a;dt)$ by Fourier
orthogonality over periods of length $2\pi$.
The missing intervals $t<1/|k|$ are controlled by
the matrix $1/\max(k,\ell)$ on $\ell^2$ for positive
integers, applied to coefficients $\sqrt{k}|c_k^\pm|$.
Its boundedness follows from Schur's test with
weight $k^{-1/2}$. The same matrix controls
the large-argument $O(z^{-1/2})$ remainders
because $\int_{\max(1/k,1/\ell)}^a t^{-2}\,dt
\le\min(k,\ell)$. Finally, the subtraction of $v_k$
on $t\ge1/|k|$ costs at most
$C|B_k|\sqrt{\log(2+|k|)}$ per mode in $L^2(dt/t)$.
Its sum is bounded already for $\alpha>1$.
These estimates give uniform pointwise-in-$\theta$
bounds on $G,E$ in $L^2(dt/t)$.
By approximation from finite Fourier data,
the limiting formula and its local Lipschitz bound
follow for $\alpha>3/2$.

For sharpness, linearize \eqref{eq:modulus-identity}
at a homogeneous field with nonzero velocity in at
least one component, using an even perturbation in
that component and fixed metric anchor.
The quadratic integral contributes zero at first order.
The derivative of the first line reduces to
$2v\cdot(\delta U_\infty-\delta U(a))$ at the chosen
angle. Take $\delta U(a)=0$.
On the resonant multiples of the rational
denominator, Cramer's rule gives
\[
 B_k=J_0(ka)\widehat{\delta U}(b,k)/
       [J_0(ka)Y_0(kb)-Y_0(ka)J_0(kb)] .
\]
On a positive-density subsequence the asymptotic cosine
of $J_0(ka)$ has magnitude bounded below;
this follows from periodicity or equidistribution
of the rotation. On these modes the limiting
intercept has multiplier of order $k^{3/2}\log k$.
Translation of the even perturbations to the
measurement angle preserves the argument.
They satisfy the periodic momentum constraint since
their spatial derivative is odd about that angle
and their initial position was fixed and homogeneous.
A bounded derivative acting on the snapshot
$\mathcal H^2_\alpha$ norm would require convergence
of a sum comparable to
$\sum k^{-1}(\log k)^{2-2\alpha}$ on the indicated
subsequence. This fails for $\alpha\le3/2$.
Taking finite real Fourier perturbations and
differentiating any proposed local Lipschitz estimate
shows its failure on smooth admissible data.
\end{proof}

\subsection{Averaging the metric datum changes its critical norm}
\label{sec:modulus-aperture}

Let $\chi\in L^\infty(\mathbb T)$ be a nonnegative
normalized spatial weight of positive measure, and
write $\langle\chi,f\rangle=\int_{\mathbb T}\chi f$.
Supply the averaged anchor $\langle\chi,\lambda(\tau_a)\rangle$.
The corresponding absolute metric datum is the limit
$\langle\chi,\lambda_\infty\rangle$.
Averaging \eqref{eq:modulus-identity} determines it
from the same field snapshots.

\begin{proposition}
\label{prop:averaged-modulus}
For the fixed rational observation gap of
Theorem~\ref{th:metric-modulus}, the averaged modulus
is locally Lipschitz in the unweighted $H^2$ snapshot
norm. For every $r<2$ it fails to be continuous
near homogeneous data in the $H^r$ snapshot norm.
If the two fields instead use stable frequency-assigned
displacement measurements in the usual wave-energy
normalization, the corresponding sharp averaged
threshold is $H^1$. In that design the sharp
pointwise modulus threshold is
$\mathcal H^1_\alpha$ with $\alpha>3/2$,
where the Fourier weight is
$\langle k\rangle[\log(2+|k|)]^\alpha$.
\end{proposition}

\begin{proof}
For a nonzero Bessel mode, set $z=|k|t$ and write
$E_k=t\,\widehat U_\theta$. The large-argument bounds
give
\[
 \int_{\tau_a}^\infty
      (|(U_\tau-v)_k|^2+|E_k|^2)\,d\tau
 \le C\langle k\rangle(|A_k|^2+|B_k|^2).
\]
One may integrate the leading $O(z)$ energy
over $1\le z\le |k|a$, the constant subtraction
over the same interval, and the integrable
small-argument remainders over $z<1$.
Parseval therefore controls the time integral in
the averaged identity at any bounded $\chi$
by the coefficient norm with weight
$\langle k\rangle^{1/2}$.
The products involving $v,U_\infty$ are bounded
in the same coefficient norm by the spatial
$L^2$ Cauchy--Schwarz inequality, since
$(\log|k|)^2=O(\langle k\rangle)$.
For a fixed rational two-snapshot gap this
coefficient norm is bounded by the snapshot $H^2$
norm. The averaging functional is quadratic plus
the supplied anchor, so the same estimates give
local Lipschitz continuity.

For sharpness, choose a single even perturbation
of a homogeneous field with zero datum at $a$.
Let its only later datum have amplitude $k^{-2}$
on a resonant multiple of the denominator.
Its Bessel coefficients have squared norm
of order $k^{-1}$, since the Bessel determinant
is of order $k^{-2}$ and
$J_0(ka)^2+Y_0(ka)^2\asymp k^{-1}$.
In the mean, the principal large-argument
energy has integrand
$(2/\pi)z(|A_k|^2+|B_k|^2)$ up to the fixed
real Fourier normalization.
This follows by adding the two squared
quadrature phases of the temporal and angular
terms. Its time integral is of order
$k(|A_k|^2+|B_k|^2)$, whereas the subtracted
limiting energy and the lower-order terms
contribute at most
$O(\log k)(|A_k|^2+|B_k|^2)$.
For $\chi$ the same leading contribution
survives: $\cos^2(k\theta)$ and
$\sin^2(k\theta)$ each have averaged value
$1/2+o(1)$ by the Riemann--Lebesgue lemma.
The linear term about the homogeneous
field is of size
$O(k^{-1/2}\log k)$ even without angular
averaging. Thus the modulus changes by a
quantity bounded away from zero while the
$H^r$ data tend to zero for $r<2$.

Under a stable frequency-assigned scheme, the
one-derivative resonant loss in the fixed-pair
inverse is absent. The Bessel coefficient
weight $\langle k\rangle^{1/2}$ is then
bounded by the assigned displacements in $H^1$.
The same proof with amplitude $k^{-1}$ for the
later data gives failure below $H^1$.
The pointwise proof of
Theorem~\ref{th:metric-modulus} uses the same
coefficients with an extra logarithmic weight, so
its upper threshold shifts down by one derivative.
For its lower bound the normalized scalar gains
are bounded above as well as below. Along a
positive-density set of modes with
$|J_0(ka)|\gtrsim k^{-1/2}$ the intercept
multiplier has order $k^{1/2}\log k$.
The same logarithmic sum gives the claimed
sharp endpoint.
\end{proof}

For a fixed smooth aperture profile $\chi_1$ of mass
one, its periodized width-$\delta$ rescaling
$\chi_\delta$ has rapidly decreasing Fourier
weights past frequency $1/\delta$.
At the homogeneous background of
Theorem~\ref{th:metric-modulus}, the $H^2$ norm of the
linearized fixed-pair averaged-modulus functional
satisfies
\[
 \|D\langle\chi_\delta,\lambda_\infty\rangle\|
      \asymp |\log\delta|^{3/2}
      \quad\text{as }\delta\downarrow0 .
\]
Indeed, on the indicated resonant subsequence the
squared norm is comparable to the partial sum
$\sum_{k\lesssim1/\delta}(\log k)^2/k$.
The aperture multiplier stays bounded below on a
fixed inner part of that range since its Fourier
transform equals one at zero. Its rapid decay
controls the remaining modes. This describes the
linearized cost of approaching the pointwise
metric convention with a narrowing spatial aperture.

\paragraph{An exact expression for the circle mean.}
For normalized integration over the circle write
\[
 E=\overline{|U_\tau|^2+e^{-2\tau}|U_\theta|^2},
 \qquad G=\overline{U\cdot U_\tau},
 \qquad \bar\lambda=\overline{\lambda}.
\]
In the polarized equations, periodic integration by parts yields
\[
 E'=-2e^{-2\tau}\overline{|U_\theta|^2},\qquad
 G'=\overline{|U_\tau|^2-e^{-2\tau}|U_\theta|^2}.
\]
Since $\bar\lambda'=E$, the combination
$\bar\lambda+E-G$ is constant. Evaluating it at
$\tau_a$ and in the asymptotic expansion gives the
quadrature-free identity
\begin{equation}
 \overline{\lambda_\infty}
 =\bar\lambda(\tau_a)+E(\tau_a)-G(\tau_a)
       -\overline{|v|^2}+\overline{U_\infty\cdot v}.
 \label{eq:mean-modulus-virial}
\end{equation}
For full vacuum wave maps, use
$G=\overline{P P_\tau}$ and the same wave-map energy
$E=\overline{|U_\tau|_{\mathbb H^2}^2+
                 e^{-2\tau}|U_\theta|_{\mathbb H^2}^2}$.
The exact derivative is instead
\[
 (\bar\lambda+E-G)'=
 \overline{(1-P)e^{2P}(Q_\tau^2-e^{-2\tau}Q_\theta^2)} .
\]
This identity holds at regular times. In an AVTD class
where the stated limits of $E,G$ and $\bar\lambda$ exist
and the right-hand side is integrable, it gives the
corresponding mean limit by adding its time integral.
The polarized identity requires no such measurement of
the transverse source.
